\paragraph{Selecting a program from the question.}
Using the $100$-task development set, we fit $\pi_Z$ to predict the member
with the highest available-repeat mean correctness on each task;
training-label ties follow member-ID order. Five task folds keep each
task's repeats together. The fitted models share a TF-IDF vocabulary
learned from development questions and use training-fold member
accuracies as additional inputs. At inference, we average the fitted
fold models' probability vectors and select the top member, falling
back to \bare{} when the top-two gap is below $0.15$, including a tie
for the highest probability. The policy is frozen before evaluation;
Appendix~\ref{app:selection} specifies fitting and fallback details.
Comparators include the development-selected fixed member and one
randomly drawn fixed member.

\paragraph{Allocating executions to a portfolio.}
An offline oracle replay compares a development-ranked portfolio with
same-code executions at budgets $b\in\{3,9,15,27\}$. Each included
member contributes three repeats. For clones, we compute exact mean
coverage over uniform size-$b$ subsets of the $27$ recorded executions.
Both arms use post-execution correctness to measure available answers.
The recorded selector choices and portfolio order stay fixed for the
recovery analysis. Replay uses tasks complete in both arms; its budget
unit is a harness execution, with model-call and token use determined
by the program. Appendix~\ref{app:selection} gives the policy, replay
formula, and accounting details.
